%%%%%%%%%%%%%%%%%%%%%%%%%%%%%%%%%%%%%%%%%%%%
%%%%%%%%%%%%%%%%%%%%%%%%%%%%%%%%%%%%%%%%%%%%%
\subsection{Forecasting Task on the Sepsis Subset of PMOA-TTS}

\begin{itemize}
    \item \textcolor{red}{We also did forecasting for additional open source and medically fine-tuned models}
    \item \textcolor{red}{Big table and detailed description of methods (task description, modeling) and results take up too much space. Push most methods in Appendix?. Only describe summarized models in main?}
    \item \textcolor{red}{Add computational resources for forecasting?}
\end{itemize}


To illustrate the potential of PMOA–TTS for downstream forecasting applications, \textcolor{red}{we include previously reported results from a small subset of sepsis case reports \citep{noroozizadeh2025reconstructing}. [Cite T2S2 or not?]}. In that study, LLM-extracted clinical textual time series (generated by Llama-3.3-70B) were used to train forecasting models for outcome prediction. While we do not rerun the experiment here, the relevant sepsis subset is now part of the publicly released PMOA–TTS corpus. This inclusion enables reproducibility of earlier findings and highlights how similar forecasting tasks can now be scaled across the full dataset or to other disease-specific cohorts. 
%In parallel work by \cite{noroozizadeh2025forecasting}, the utility of textual time series generated by Llama-3.3-70B for forecasting tasks has been demonstrated on a smaller cohort of sepsis patients. 
We define the following two main forecasting tasks.

\paragraph{Event Occurrence Prediction:} In this task, the model receives a prefix of the clinical timeline—comprising all events up to a given time point $t$—and must predict whether each of the next $k$ events will occur within a defined time window. This simulates an “online” prediction scenario, where, for each of the next $k$ events, the model outputs a binary label: does this event happen within $h$ hours after $t$? Time horizons considered include 1 hour, 24 hours (1 day), and 168 hours (1 week). The task is structured as a sequence of binary classification problems, with evaluation based on precision, recall, and F1 score, averaged across the $k$ events for each time window.

\paragraph{Temporal Ordering Prediction:} This task evaluates the model’s ability to infer the correct chronological sequence of future events. At each cutoff time $t$, we select the subsequent $k$ events and remove their timestamps. The model is then required to output a permutation of these events that aligns with their true temporal order. This is formulated as a ranking problem, and performance is measured by the pairwise concordance—i.e., the proportion of correctly ordered event pairs—between the predicted and actual sequences. This setup challenges the model to infer temporal flow based solely on event content.


%%%%%%%%%%%%%%%%%%%%%%%%%%%%%%%%%%%%%%%%%%%%%%%%
\subsubsection{Modeling approaches}

To evaluate performance on the event forecasting tasks, we use three modeling paradigms: (i) fine-tuned large language models (LLMs) with task-specific heads, (ii) prompted LLMs in zero- or few-shot settings without gradient updates, and (iii) BERT-style encoders with fine-tuned heads tailored to each task. 

\paragraph{Fine-tuned LLMs:} We use instruction-tuned decoder-only models from the \texttt{LLaMA} and \texttt{DeepSeek} families, including \texttt{Llama-3.3-70B-Instruct}, \texttt{Llama-3.1-8B-Instruct}, and their distilled variants. Each is paired with a lightweight MLP head for classification or ranking, depending on the task. Inputs include a text-formatted event prefix (e.g., timeline up to time $t$), optionally with an instruction template. The output layer generates task-specific predictions—binary labels for event occurrence or a permutation over $k$ events for ordering—trained using cross-entropy or pairwise ranking loss.

\paragraph{Prompted LLMs:} In the zero- or few-shot setting, we use the same LLM architectures without fine-tuning, guiding outputs via structured prompts at inference. Each prompt includes: (1) a system role (e.g., "You are an expert physician."), (2) a task-specific user instruction, and (3) one or more few-shot examples. Prompts are tailored per task and described in \cite{noroozizadeh2025forecasting}. Model outputs are parsed into binary labels or ordered lists; ambiguous or malformed generations are excluded from evaluation.

\paragraph{Fine-tuned Encoder-Only Models:} We also evaluate encoder-based models trained end-to-end for each task, including \texttt{BERT-base-uncased}, \texttt{RoBERTa-base}, \texttt{DeBERTa-v3-small}, \texttt{ModernBERT-small}, and \texttt{ModernBERT-large}. Each is paired with a task-specific MLP head for event occurrence or ordering prediction. Inputs use the same tokenized event sequence prefixes as in other settings, formatted per model requirements. Models are trained with standard supervised objectives and evaluated using F1 score for event occurrence and pairwise concordance for ordering.


% \begin{table}[htb!]
\scriptsize
\caption{Forecasting performance (event occurrence: F1 and correct event ordering: concordance-index) of the ensuing $k=8$ events. All models are trained/fine-tuned on time-ordered annotations from \texttt{DeepSeek-R1}. \textbf{Bold} values indicate best in category amongst all models.}
\label{tab:DS_forecasting_clean}
\centering
\resizebox{\linewidth}{!}{
\begin{tabular}{lcccccccc}
\toprule
& \multicolumn{4}{c}{T2S2 (DeepSeek R1)} & \multicolumn{2}{c}{sepsis-10} & \multicolumn{2}{c}{sepsis-100}\\
& F1(1h) & F1(1d) & F1(1w) & c-index & F1(1d) & c-index & F1(1d) & c-index\\
\midrule
\textbf{LLM-MLP head} &&&&&&&&\\
DS-L3.3 70B & \textbf{0.075} & 0.482 & \textbf{0.796} & \textbf{0.632} & \textbf{0.397} & \textbf{0.578} & \textbf{0.613} & 0.595 \\
\textbf{Prompting} &&&&&&&&\\
L3.3 70B few-shot & \textbf{0.095} & \textbf{0.380} & 0.652 & -- & \textbf{0.360} & -- & 0.452 & -- \\
\midrule
\textbf{Encoder-MLP head} &&&&&&&&\\
ModernBERT-base & \textbf{0.306} & \textbf{0.576} & \textbf{0.877} & \textbf{0.646} & 0.428 & \textbf{0.562} & 0.607 & \textbf{0.558} \\
\midrule
\textbf{Encoder-masking-fine-tuned} &&&&&&&&\\
ModernBERT-base & \textbf{0.186} & 0.449 & \textbf{0.697} & \textbf{0.672} & 0.325 & 0.576 & \textbf{0.503} & 0.613 \\
\midrule
\textbf{Encoder-masking-zeroshot} &&&&&&&&\\
ModernBERT-base & 0.169 & \textbf{0.468} & \textbf{0.804} & 0.496 & \textbf{0.358} & 0.501 & \textbf{0.511} & 0.499 \\
ModernBERT-large & 0.133 & 0.349 & 0.534 & 0.494 & 0.271 & 0.509 & 0.277 & 0.502 \\
\bottomrule
\end{tabular}
}
\end{table}

%\begin{table}[tb!]
\scriptsize
\caption{Forecasting performance for textual time series for sepsis (T2S2) (event occurrence: F1 and correct event ordering: concordance-index) of the ensuing $k=8$ events. All models are trained/fine-tuned on time-ordered annotations from \texttt{Llama 3.3-70B}. \textbf{Bold} values indicate best in category amongst all models. The sepsis-10 and sepsis-100 sets are evaluation subsets, consisting of 10 expert-annotated case reports serving as the gold standard 100 sepsis reports for broader model evaluation respectively.}
\label{tab:L33_forecasting_clean}
\centering
\resizebox{\linewidth}{!}{
\begin{tabular}{lcccccccc}
\toprule
& \multicolumn{4}{c}{T2S2 (Llama 3.3)} & \multicolumn{2}{c}{sepsis-10} & \multicolumn{2}{c}{sepsis-100} \\
& F1(1h) & F1(1d) & F1(1w) & c-index & F1(1d) & c-index & F1(1d) & c-index \\
\midrule
\textbf{LLM-MLP head} &&&&&&&&\\
DS-L3.3 70B & 0.140 & 0.563 & 0.811 & \textbf{0.624} & \textbf{0.397} & \textbf{0.585} & \textbf{0.629} & \textbf{0.598} \\
\textbf{Prompting} &&&&&&&&\\
L3.3 70B few-shot ordinal & 0.064 & \textbf{0.460} & \textbf{0.729} & -- & 0.313 & -- & 0.480 & -- \\
\midrule
\textbf{Encoder-MLP head} &&&&&&&&\\
ModernBERT-base & 0.257 & \textbf{0.645} & \textbf{0.903} & 0.594 & 0.431 & \textbf{0.590} & \textbf{0.626} & 0.565 \\
\midrule
\textbf{Encoder-masking-fine-tuned} &&&&&&&&\\
ModernBERT-base & \textbf{0.139} & \textbf{0.489} & 0.700 & 0.632 & \textbf{0.334} & 0.595 & \textbf{0.478} & 0.612 \\
\midrule
\textbf{Encoder-masking-zeroshot} &&&&&&&&\\
ModernBERT-base & \textbf{0.132} & \textbf{0.525} & \textbf{0.840} & 0.556 & \textbf{0.352} & 0.503 & 0.511 & 0.501 \\
\bottomrule
\end{tabular}
}
\end{table}

% \begin{table}[htb!]
\scriptsize
\label{L33_forecasting}
\centering
\resizebox{\linewidth}{!}{
\begin{tabular}{lcccccccc}
\toprule
& \multicolumn{4}{c}{T2S2 (Llama 3.3)} & \multicolumn{2}{c}{sepsis-10} & \multicolumn{2}{c}{sepsis-100} \\
\cmidrule(lr){2-5} \cmidrule(lr){6-7} \cmidrule(lr){8-9}
& F1(1h) & F1(1d) & F1(1w) & c-index & F1(1d) & c-index & F1(1d) & c-index \\
\midrule
\textbf{LLM-MLP head} &&&&&&&&\\
DS-L3.3 70B & 0.140 & 0.563 & 0.811 & \textbf{0.624} & \textbf{0.397} & \textbf{0.585} & \textbf{0.629} & \textbf{0.598} \\
L3.3 70B & 0.265 & 0.466 & 0.807 & 0.619 & 0.160 & 0.567 & 0.466 & 0.597 \\
DS-L3.1 8B & 0.264 & \textbf{0.584} & \textbf{0.837} & \textbf{0.624} & 0.258 & 0.576 & 0.605 & \textbf{0.598} \\
L3.1 8B & \textbf{0.338} & 0.565 & 0.823 & 0.619 & 0.368 & 0.569 & 0.536 & 0.596 \\
\textbf{Prompting} &&&&&&&&\\
L3.3 70B 0-shot & 0.085 & 0.443 & 0.616 & -- & 0.236 & -- & 0.452 & -- \\
L3.3 70B 0-shot ordinal & 0.063 & 0.454 & 0.707 & -- & 0.289 & -- & \textbf{0.483} & -- \\
L3.3 70B few-shot & \textbf{0.099} & 0.457 & 0.703 & -- & \textbf{0.360} & -- & 0.452 & -- \\
L3.3 70B few-shot ordinal & 0.064 & \textbf{0.460} & \textbf{0.729} & -- & 0.313 & -- & 0.480 & -- \\
\midrule
\textbf{Encoder-MLP head} &&&&&&&&\\
BERT & 0.224 & 0.622 & 0.899 & 0.500 & 0.400 & 0.500 & 0.612 & 0.500 \\
RoBERTa & 0.180 & 0.608 & 0.900 & 0.500 & 0.364 & 0.500 & 0.597 & 0.500 \\
DeBERTa-small & 0.258 & 0.641 & 0.898 & \textbf{0.603} & 0.427 & 0.585 & 0.616 & \textbf{0.592} \\
ModernBERT-base & 0.257 & \textbf{0.645} & \textbf{0.903} & 0.594 & \textbf{0.431} & \textbf{0.590} & \textbf{0.626} & 0.565 \\
ModernBERT-large & \textbf{0.312} & 0.622 & 0.897 & 0.500 & 0.404 & 0.500 & 0.623 & 0.500 \\
\midrule
\textbf{Encoder-masking-fine-tuned} &&&&&&&&\\
BERT & 0.131 & 0.471 & \textbf{0.741} & 0.622 & 0.320 & 0.580 & 0.467 & 0.622 \\
RoBERTa & 0.129 & 0.484 & 0.659 & 0.625 & 0.329 & 0.587 & 0.472 & 0.614 \\
DeBERTa-small & 0.135 & 0.468 & 0.693 & \textbf{0.648} & 0.318 & \textbf{0.602} & 0.462 & 0.623 \\
ModernBERT-base & \textbf{0.139} & \textbf{0.489} & 0.700 & 0.632 & \textbf{0.334} & 0.595 & \textbf{0.478} & 0.612 \\
ModernBERT-large & 0.124 & 0.470 & 0.712 & 0.646 & 0.318 & 0.586 & 0.468 & \textbf{0.631} \\
\midrule
\textbf{Encoder-masking-zeroshot} &&&&&&&&\\
BERT & 0.000 & 0.000 & 0.000 & 0.504 & 0.000 & 0.515 & 0.000 & 0.506 \\
RoBERTa & 0.074 & 0.279 & 0.445 & 0.500 & 0.083 & 0.500 & 0.319 & 0.500 \\
DeBERTa-small & 0.008 & 0.439 & 0.839 & \textbf{0.562} & 0.296 & \textbf{0.518} & \textbf{0.518} & 0.501 \\
ModernBERT-base & \textbf{0.132} & \textbf{0.525} & \textbf{0.840} & 0.556 & \textbf{0.352} & 0.503 & 0.511 & 0.501 \\
ModernBERT-large & 0.109 & 0.413 & 0.622 & 0.499 & 0.326 & 0.509 & 0.339 & \textbf{0.518} \\
\bottomrule
\end{tabular}
}
\caption{Forecasting performance (event occurrence: F1 and correct event ordering: concordance-index) of the ensuing $k=8$ events on the sepsis subset of PMOA-TTS (T2S2:\textbf{T}extual \textbf{T}ime-\textbf{S}eries for \textbf{S}epsis). \textbf{All models are trained/fine-tuned on time-ordered annotations from \texttt{Llama-3.3-70B}.} The sepsis-10 and sepsis-100 sets are evaluation subsets, consisting of 10 expert-annotated case reports serving as the gold standard and 100 sepsis reports for broader model evaluation respectively. \textcolor{red}{Need to add new results for open-source and medically fine-tuned models.}}
\label{L33_forecasting}
\vspace{-2em}
\end{table}

\subsection{Forecasting results}

We show the forecasting results for 2319 textual time series for sepsis (T2S2). \textcolor{red}{Cite T2S2 or not?}. From this set, we randomly selected two subsets for evaluation: a group of 10 reports (sepsis-10), which underwent expert clinical annotation to serve as a gold standard, and a larger group of 100 reports (sepsis-100) used for broader testing. 

\begin{table}[htb!]
\scriptsize
\label{L33_forecasting}
\centering
\resizebox{\linewidth}{!}{
\begin{tabular}{lcccccccc}
\toprule
& \multicolumn{4}{c}{T2S2 (Llama 3.3)} & \multicolumn{2}{c}{sepsis-10} & \multicolumn{2}{c}{sepsis-100} \\
\cmidrule(lr){2-5} \cmidrule(lr){6-7} \cmidrule(lr){8-9}
& F1(1h) & F1(1d) & F1(1w) & c-index & F1(1d) & c-index & F1(1d) & c-index \\
\midrule
\textbf{LLM-MLP head} &&&&&&&&\\
DS-L3.3 70B & 0.140 & 0.563 & 0.811 & \textbf{0.624} & \textbf{0.397} & \textbf{0.585} & \textbf{0.629} & \textbf{0.598} \\
L3.3 70B & 0.265 & 0.466 & 0.807 & 0.619 & 0.160 & 0.567 & 0.466 & 0.597 \\
DS-L3.1 8B & 0.264 & \textbf{0.584} & \textbf{0.837} & \textbf{0.624} & 0.258 & 0.576 & 0.605 & \textbf{0.598} \\
L3.1 8B & \textbf{0.338} & 0.565 & 0.823 & 0.619 & 0.368 & 0.569 & 0.536 & 0.596 \\
\textbf{Prompting} &&&&&&&&\\
L3.3 70B 0-shot & 0.085 & 0.443 & 0.616 & -- & 0.236 & -- & 0.452 & -- \\
L3.3 70B 0-shot ordinal & 0.063 & 0.454 & 0.707 & -- & 0.289 & -- & \textbf{0.483} & -- \\
L3.3 70B few-shot & \textbf{0.099} & 0.457 & 0.703 & -- & \textbf{0.360} & -- & 0.452 & -- \\
L3.3 70B few-shot ordinal & 0.064 & \textbf{0.460} & \textbf{0.729} & -- & 0.313 & -- & 0.480 & -- \\
\midrule
\textbf{Encoder-MLP head} &&&&&&&&\\
BERT & 0.224 & 0.622 & 0.899 & 0.500 & 0.400 & 0.500 & 0.612 & 0.500 \\
RoBERTa & 0.180 & 0.608 & 0.900 & 0.500 & 0.364 & 0.500 & 0.597 & 0.500 \\
DeBERTa-small & 0.258 & 0.641 & 0.898 & \textbf{0.603} & 0.427 & 0.585 & 0.616 & \textbf{0.592} \\
ModernBERT-base & 0.257 & \textbf{0.645} & \textbf{0.903} & 0.594 & \textbf{0.431} & \textbf{0.590} & \textbf{0.626} & 0.565 \\
ModernBERT-large & \textbf{0.312} & 0.622 & 0.897 & 0.500 & 0.404 & 0.500 & 0.623 & 0.500 \\
\midrule
\textbf{Encoder-masking-fine-tuned} &&&&&&&&\\
BERT & 0.131 & 0.471 & \textbf{0.741} & 0.622 & 0.320 & 0.580 & 0.467 & 0.622 \\
RoBERTa & 0.129 & 0.484 & 0.659 & 0.625 & 0.329 & 0.587 & 0.472 & 0.614 \\
DeBERTa-small & 0.135 & 0.468 & 0.693 & \textbf{0.648} & 0.318 & \textbf{0.602} & 0.462 & 0.623 \\
ModernBERT-base & \textbf{0.139} & \textbf{0.489} & 0.700 & 0.632 & \textbf{0.334} & 0.595 & \textbf{0.478} & 0.612 \\
ModernBERT-large & 0.124 & 0.470 & 0.712 & 0.646 & 0.318 & 0.586 & 0.468 & \textbf{0.631} \\
\midrule
\textbf{Encoder-masking-zeroshot} &&&&&&&&\\
BERT & 0.000 & 0.000 & 0.000 & 0.504 & 0.000 & 0.515 & 0.000 & 0.506 \\
RoBERTa & 0.074 & 0.279 & 0.445 & 0.500 & 0.083 & 0.500 & 0.319 & 0.500 \\
DeBERTa-small & 0.008 & 0.439 & 0.839 & \textbf{0.562} & 0.296 & \textbf{0.518} & \textbf{0.518} & 0.501 \\
ModernBERT-base & \textbf{0.132} & \textbf{0.525} & \textbf{0.840} & 0.556 & \textbf{0.352} & 0.503 & 0.511 & 0.501 \\
ModernBERT-large & 0.109 & 0.413 & 0.622 & 0.499 & 0.326 & 0.509 & 0.339 & \textbf{0.518} \\
\bottomrule
\end{tabular}
}
\caption{Forecasting performance (event occurrence: F1 and correct event ordering: concordance-index) of the ensuing $k=8$ events on the sepsis subset of PMOA-TTS (T2S2:\textbf{T}extual \textbf{T}ime-\textbf{S}eries for \textbf{S}epsis). \textbf{All models are trained/fine-tuned on time-ordered annotations from \texttt{Llama-3.3-70B}.} The sepsis-10 and sepsis-100 sets are evaluation subsets, consisting of 10 expert-annotated case reports serving as the gold standard and 100 sepsis reports for broader model evaluation respectively. \textcolor{red}{Need to add new results for open-source and medically fine-tuned models.}}
\label{L33_forecasting}
\vspace{-2em}
\end{table}


\subsubsection{Event forecast within next 24 hours: F1 performance} Our results show that encoder-based models outperform LLMs in event forecasting across \texttt{Llama-3.3-70B} (Table ~\ref{L33_forecasting}) annotations, especially in short- and medium-term prediction windows. Fine-tuned encoder models with MLP heads achieve significantly higher F1 scores than LLMs, confirming the strength of encoder-based representations for forecasting. Among them, MLP fine-tuning consistently outperforms masking-based methods, particularly for long-horizon forecasting. \texttt{ModernBERT-base} achieves the highest F1 across all horizons, especially at 168 hours, with \texttt{ModernBERT-large} performing similarly but without consistent gains. Zero-shot masking models like \texttt{BERT} and \texttt{RoBERTa} perform poorly, while \texttt{DeBERTa-small} and \texttt{ModernBERT-base} show stronger generalization.

Performance improves with longer forecasting windows, indicating better predictability of long-term trends and greater difficulty in short-term forecasting due to event sparsity. Models also generalize poorly to external datasets (sepsis-10, sepsis-100), especially at the 1-day horizon, highlighting challenges in domain adaptation. Nonetheless, \texttt{ModernBERT} maintains strong relative performance, demonstrating robustness across datasets.

\subsubsection{Temporal ordering of forecasted events: Concordance} In addition to event occurrence prediction, models are evaluated on their ability to rank upcoming events using concordance (c-index), based on annotations from \texttt{Llama-3.3-70B} (Table ~\ref{L33_forecasting}). Encoder models again outperform LLMs, which fail to produce reliable orderings even with fine-tuning or prompting. Within encoder architectures, training strategies impact performance: MLP-head fine-tuned models lead in F1 but do not consistently achieve the best concordance. Instead, fine-tuned masking models—particularly \texttt{ModernBERT-large}—achieve the highest c-index, indicating a strength in capturing event order over binary classification.

Dataset comparisons reveal generalization challenges. While fine-tuned models perform well on the internal T2S2 set, c-index scores drop on external datasets, especially for non-\texttt{ModernBERT} models. \texttt{ModernBERT-large} (fine-tuned masking) maintains top c-index scores across all datasets, confirming its robustness. MLP-head models retain strong F1 scores but slightly lag in concordance, especially at the 168-hour horizon, suggesting trade-offs between classification and ranking.

Overall, \texttt{ModernBERT} models are the strongest across F1 and concordance, with MLP-head fine-tuning excelling at binary tasks and masking fine-tuning better suited for event ordering. However, all models struggle with external generalization, underscoring the need for improved domain adaptation.